% =====================================================================
% Приложение А. Функциональная схема программной системы.
% Реализовано в TikZ. Версия 4 — горизонтальный snake-flow на
% альбомной странице, слои представлены крупными блоками с перечнем
% содержимого внутри, без \resizebox.
% =====================================================================

\clearpage
\section{Функциональная схема системы}
\label{app:uml}

\renewcommand{\thefigure}{\thesection.\arabic{figure}}
\setcounter{figure}{0}

В настоящем приложении приведена функциональная схема программной
системы, описанной в главе~\ref{sec:architecture}. Схема изображает
движение управления и данных между основными архитектурными блоками
от поступления входных данных эксперимента до получения результирующих
метрик и журналов. Поток выполнения организован по~двум рядам слева
направо, переход на~второй ряд показан изгибом стрелки справа.
На~рисунке~\ref{fig:appA:scheme} двойная красная рамка отмечает
авторские компоненты~--- слой адаптации, оракул-маршрутизатор и~шлюз
взаимодействия с~человеком.

\begin{sidewaysfigure}[p]
\centering
\begin{tikzpicture}[
  font=\small,
  every node/.style={align=center},
  io/.style={rectangle, draw=black, thick, fill=gray!12,
             text width=3.4cm, minimum height=2.0cm, rounded corners=3pt},
  ctrl/.style={rectangle, draw=black, thick, fill=blue!10,
               text width=3.4cm, minimum height=2.0cm, rounded corners=3pt},
  authctrl/.style={rectangle, draw=red!75!black, very thick, double,
                   double distance=1.5pt, fill=blue!10,
                   text width=3.4cm, minimum height=2.0cm, rounded corners=3pt},
  topo/.style={rectangle, draw=black, thick, fill=orange!20,
               text width=3.4cm, minimum height=2.0cm, rounded corners=3pt},
  agent/.style={rectangle, draw=black, thick, fill=green!18,
                text width=3.4cm, minimum height=2.0cm, rounded corners=3pt},
  infra/.style={rectangle, draw=black, thick, fill=red!12,
                text width=3.4cm, minimum height=2.0cm, rounded corners=3pt},
  authinfra/.style={rectangle, draw=red!75!black, very thick, double,
                    double distance=1.5pt, fill=red!12,
                    text width=3.4cm, minimum height=2.0cm, rounded corners=3pt},
  eval/.style={rectangle, draw=black, thick, fill=violet!15,
               text width=3.4cm, minimum height=2.0cm, rounded corners=3pt},
  autheval/.style={rectangle, draw=red!75!black, very thick, double,
                   double distance=1.5pt, fill=violet!15,
                   text width=3.4cm, minimum height=2.0cm, rounded corners=3pt},
  store/.style={rectangle, draw=black, thick, fill=yellow!22,
                text width=3.4cm, minimum height=2.0cm, rounded corners=3pt},
  flow/.style={-{Latex[length=3mm]}, very thick},
]

% ===== Ряд 1 (сверху) =====
\node[io]            (input)    at (0,   0)  {\textbf{Вход}\\Задача $\mathcal{T}$\\и YAML-конфиг};
\node[ctrl]          (orch)     at (4.0, 0)  {\textbf{Оркестратор}\\\code{atm.experiment}\\инициализация};
\node[authctrl]      (adapt)    at (8.0, 0)  {\textbf{Слой адаптации $\bigstar$}\\PhaseRouter\\TopologyRouter\\SwitchGuards};
\node[topo]          (topo)     at (12.0, 0) {\textbf{Топологии}\\Star, Chain, Mesh,\\Debate, Hierarchical\\(активна одна)};
\node[agent]         (agents)   at (16.0, 0) {\textbf{Агенты}\\Planner, Researcher,\\Executor, Critic,\\Debater, Coordinator};
\node[infra]         (services) at (20.0, 0) {\textbf{Внешние службы}\\Инструменты+Docker\\\code{LLMWrapper}};

% ===== Ряд 2 (снизу, движение справа налево) =====
\node[authinfra]     (human)    at (20.0, -5.0) {\textbf{Шлюз человека $\bigstar$}\\\code{HumanGateway}\\(пять ролей)};
\node[ctrl]          (obs)      at (16.0, -5.0) {\textbf{Наблюдаемость}\\\code{atm.observability}\\захват событий};
\node[eval]          (judge)    at (12.0, -5.0) {\textbf{Слой оценки}\\\code{atm.evaluation}\\судья + NASA-TLX};
\node[autheval]      (oracle)   at (8.0, -5.0)  {\textbf{Оракул $\bigstar$}\\per-task best-of\\$\Delta_{\text{oracle}}$};
\node[store]         (store)    at (4.0, -5.0)  {\textbf{Хранение}\\PostgreSQL 16\\Apache Parquet};
\node[io]            (output)   at (0,   -5.0)  {\textbf{Выход}\\Журналы +\\метрики $q, c, \tau,$\\$N_{\text{switch}}, r_{\text{guard}}$};

% ===== Связи ряда 1 =====
\draw[flow] (input)    -- (orch);
\draw[flow] (orch)     -- (adapt);
\draw[flow] (adapt)    -- (topo);
\draw[flow] (topo)     -- (agents);
\draw[flow] (agents)   -- (services);

% ===== Изгиб с ряда 1 на ряд 2 =====
\draw[flow] (services.east) -- ++(0.6, 0)
                            |- (human.east);

% ===== Связи ряда 2 (справа налево) =====
\draw[flow] (human)    -- (obs);
\draw[flow] (obs)      -- (judge);
\draw[flow] (judge)    -- (oracle);
\draw[flow] (oracle)   -- (store);
\draw[flow] (store)    -- (output);

% ===== Обратная связь: оракул → слой адаптации =====
\draw[flow, dashed, red!70!black]
    (oracle.north) -- ++(0, 1.2)
    node[midway, right=2pt, font=\footnotesize, black]
       {оракульная таблица}
    -| (adapt.south);

% ===== Обратная связь: шлюз человека → слой адаптации =====
\draw[flow, dashed, red!70!black]
    (human.north) -- ++(0, 1.7)
    node[midway, right=2pt, font=\footnotesize, black]
       {HITL-сигналы}
    -| (adapt.south east);

\end{tikzpicture}

\caption{Функциональная схема программной системы. Поток управления
организован по~двум рядам слева направо с~переходом на~второй ряд
справа. Двойная красная рамка обозначает авторские компоненты~---
слой адаптации с~маркером $\bigstar$ (маршрутизаторы фаз
и~топологий с~защитными правилами \code{SwitchGuards}),
шлюз взаимодействия с~человеком и~оракул-маршрутизатор. Штриховые
красные стрелки --- обратные связи: оракульная таблица возвращается
в~маршрутизатор топологий, HITL-сигналы поступают от~человека
к~слою адаптации. Полный перечень заимствованных и~авторских
компонентов приведен в~таблице~\ref{tab:arch:contribution}.}
\label{fig:appA:scheme}
\end{sidewaysfigure}

\subsection*{Пояснение к схеме}

Поток выполнения начинается с поступления конфигурации эксперимента
в формате YAML и описания задачи $\mathcal{T}$ из выбранного
корпуса. Оркестратор эксперимента (\code{atm.experiment})
инициализирует общее состояние графа и передает управление слою
адаптации. Маршрутизатор фаз принимает решение о текущей фазе
решения, передает результат маршрутизатору топологий, который
выбирает активную коммуникационную топологию. В работе одновременно
действует ровно одна из пяти базовых топологий, а адаптивный режим
реализуется через последовательное переключение между ними по
решению маршрутизатора.

Активная топология определяет порядок активации агентов и протокол
обмена сообщениями между ними. Любая из шести ролей агентов
(планировщик, исследователь, исполнитель, критик, дебатер,
координатор) может быть подключена к любой топологии в соответствии
с правилами раздела~\ref{sec:meth:topologies}. Агенты обращаются
за внешними службами --- инструментами и Docker-песочницей для
исполнения кода, оболочкой над языковыми моделями для генерации
текста, шлюзом взаимодействия с человеком.

Все события системы (вызовы языковых моделей, выставление сигналов,
переходы между фазами, переключения топологий, обращения к человеку,
срабатывания защитных правил маршрутизатора) перехватываются
подсистемой наблюдаемости и направляются в два хранилища ---
реляционную базу для метаданных и колоночные файлы Apache Parquet
для объемных журналов. На финальной стадии результаты прогона
становятся доступными для постобработки слою анализа
\code{atm.analysis}.

Обратный канал от шлюза человека к маршрутизатору топологий
обозначает поступление управляющих сигналов от участника-человека
к адаптивному слою. Через этот канал человек, действующий в роли
координатора или наблюдателя, может форсировать переключение
топологии или досрочно завершить фазу проверки. Штриховая связь
от~оракула к~маршрутизатору топологий обеспечивает возврат
оракульной таблицы (per-task best-of) для расчета
$\Delta_{\text{oracle}}$.
